\documentclass[11pt]{article}
\usepackage[letterpaper,margin=1in]{geometry}
\usepackage{amsmath,amssymb,amsthm}
\usepackage{booktabs,tabularx,array}
\usepackage{needspace}
\usepackage[colorlinks=true,linkcolor=black,citecolor=black,urlcolor=blue]{hyperref}
\hypersetup{pdftitle={Three-chromatic Cayley graphs of exponent-four groups},
  pdfauthor={Alex Chan}}

\newtheorem{theorem}{Theorem}
\newtheorem{lemma}{Lemma}
\newcommand{\out}{N^{+}}
\newcommand{\Cay}{\operatorname{Cay}}
\newcommand{\F}{\mathbb F}

\title{Three-chromatic Cayley graphs of exponent-four groups}
\author{Alex Chan\\[2pt]
  \small Harvard University and National Bureau of Economic Research\\[-1pt]
  \small \texttt{achan@hbs.edu}}
\date{\today}

\begin{document}
\maketitle

\begin{abstract}
We answer Question 5.3 of Krebs and Sankar. A nonabelian group of
exponent four admits an undirected Cayley graph of chromatic number three
if and only if it is not a direct product of the quaternion group of order
eight and an elementary abelian two-group. The classification includes
infinite groups, and every positive case has a connection set of at most
ten elements.
\end{abstract}

\section{Introduction and the dihedral witness}

Krebs and Sankar classified the abelian groups admitting a three-chromatic
Cayley graph and asked for the corresponding classification among
nonabelian groups of exponent four \cite[Question~5.3]{KrebsSankar}.
The obstruction turns out to be precisely the quaternionic family.
An explicit eight-vertex graph supplies all positive cases through
sections, while a graph isomorphism reduces the negative cases to
Payan's theorem on cubelike graphs \cite{Payan}.

For a group $G$ and $S=S^{-1}\subseteq G\setminus\{1\}$, we join $x$ to
$y$ in $\Cay(G,S)$ when $x^{-1}y\in S$. In particular, the graph need not
be connected, as in the question under consideration. We work in ZFC.
A \emph{section} of $G$ means a quotient $H/N$ of a subgroup $H\leq G$.

\begin{theorem}\label{thm:main}
Let $G$ be a nonabelian group in which $g^4=1$ for all $g\in G$.
The following are equivalent:
\begin{enumerate}
\item $G$ admits a simple undirected Cayley graph of chromatic number three;
\item $G$ has a section isomorphic to the dihedral group $D_8$ of order eight;
\item $G$ is not isomorphic to $Q_8\times E$ for any elementary abelian
two-group $E$.
\end{enumerate}
The group $E$ may be infinite. In (1), the connection set can be chosen to
have at most ten elements.
\end{theorem}

Put $D_8=\langle r,s\mid r^4=s^2=1,\ srs=r^{-1}\rangle$ and
$S_0=\{r^2,s,rs\}$. These are three involutions. With vertices ordered as
$1,r,r^2,r^3,s,rs,r^2s,r^3s$, respectively, the assignment
\[
             (0,0,1,2,1,1,2,0)
\]
is a proper three-coloring of $\Cay(D_8,S_0)$. The vertices
$1,r^2,r^2s,r,rs,1$ form a five-cycle. Hence the graph has chromatic
number three; it is the eight-vertex Wagner graph.

\Needspace{10\baselineskip}
\begin{lemma}\label{lem:lift}
If $G$ has a $D_8$ section, it admits a three-chromatic Cayley graph with
at most ten elements in its connection set.
\end{lemma}
\begin{proof}
Let $\pi:H\twoheadrightarrow D_8$ be the quotient map and let
$q_0,q_1,q_2,q_3,q_4,q_5=q_0$ be the displayed five-cycle. Choose
$v_i\in\pi^{-1}(q_i)$ for $0\leq i<5$, set $v_5=v_0$, and put
\[
 S=\{v_i^{-1}v_{i+1},\ (v_i^{-1}v_{i+1})^{-1}:0\leq i<5\}.
\]
Its image under $\pi$ lies in $S_0$, so $S=S^{-1}$ and $1\notin S$.
The three-coloring of the dihedral graph pulls back to $\Cay(H,S)$,
while $v_0,v_1,v_2,v_3,v_4,v_0$ is an odd cycle. Regard $S$ as a subset
of $G$. Every component of $\Cay(G,S)$ lies in a left coset of $H$;
translate the coloring across the cosets. Thus $\chi(\Cay(G,S))=3$.
\end{proof}

\section{A structural dichotomy}

We give a direct argument without a finiteness assumption.

\begin{lemma}\label{lem:structure}
If a nonabelian group $G$ of exponent four has no $D_8$ section, then
$G\cong Q_8\times E$ for an elementary abelian two-group $E$.
\end{lemma}
\begin{proof}
Two noncommuting involutions generate $D_8$: their product has order
four and they act as two reflections. Hence all involutions of $G$
commute. The elements of square one form a normal elementary abelian
subgroup $T$.

We first show that every involution is central. Otherwise, choose
$t\in T$ and $x\in G$ such that $t'=xtx^{-1}\ne t$. Necessarily $x$
has order four. Conjugation by $x$ has order at most two on $T$, since
$x^2\in T$ and $T$ is abelian; it interchanges $t$ and $t'$. Thus
$A=\langle t,t'\rangle\cong C_2^2$ is normal in $H=\langle A,x\rangle$.
Write $w=x^2$. If $w\in A$, it is the unique nonidentity element of
$A$ fixed by the swap, namely $tt'$. Then $H$ has order eight and is
$D_8$. If $w\notin A$, the group $H$ has order sixteen and
$H/\langle w\rangle\cong A\rtimes C_2\cong D_8$. Both alternatives
contradict the hypothesis.

All squares are therefore central. Quotienting by the central subgroup
generated by squares produces a group of exponent two, hence an abelian
group. Thus $G$ has nilpotence class at most two. For all $x,y\in G$,
the commutator $[x,y]$ is central and
$[x,y]^2=[x^2,y]=1$.

Take noncommuting $x,y$ and write
$a=x^2$, $b=y^2$, $c=[x,y]$, $C=\langle a,b,c\rangle$,
$H=\langle x,y\rangle$. Here $C$ is a central elementary abelian
two-group, $c\ne1$, and $H/C\cong C_2^2$: the cosets $xC$ and $yC$
are nontrivial and distinct, since otherwise $x$ and $y$ would
commute. If $a\ne c$, there is a linear functional
$\ell:C\to C_2$ with $\ell(c)=1$ and $\ell(a)=0$; the same holds
with $b$ in place of $a$ if $b\ne c$. The quotient
$H/\ker\ell$ has order eight, nontrivial commutator, and a
noncentral involution (the image of $x$ or $y$). It is $D_8$.
Consequently every noncommuting pair satisfies
\begin{equation}\label{eq:pairs}
                     x^2=y^2=[x,y]\ne1.
\end{equation}

Choose noncommuting $i,j$ and set $z=i^2=j^2=[i,j]$.
Equation~\eqref{eq:pairs} gives
$K=\langle i,j\rangle\cong Q_8$. For every $g\in G$, each of
$[i,g]$ and $[j,g]$ is either $1$ or $z$. Conjugation by $g$ therefore
acts on $K$ by independent sign changes of $i$ and $j$; all four such
automorphisms are inner in $K$. Hence $G=K C_G(K)$.
If $h\in C_G(K)$, the pair $ih,j$ does not commute. Applying
\eqref{eq:pairs} yields $(ih)^2=z=i^2$, so $h^2=1$.
Thus $C_G(K)$ is elementary abelian (indeed central in $G$), and
$K\cap C_G(K)=\langle z\rangle$. Choose a vector-space complement
$E$ to $\langle z\rangle$ in $C_G(K)$. It follows that
$G=K\times E\cong Q_8\times E$.
\end{proof}

\Needspace{8\baselineskip}
\section{The quaternionic obstruction}

The following graph isomorphism explains the exceptional family.

\begin{lemma}\label{lem:cubelike}
Every undirected Cayley graph on $Q_8\times E$ is isomorphic to a
Cayley graph on the elementary abelian two-group $\F_2^3\times E$.
\end{lemma}
\begin{proof}
Write $z=-1\in Q_8$, and choose representatives
$t(00)=1$, $t(10)=i$, $t(01)=j$, $t(11)=k$ of
$Q_8/\langle z\rangle\cong\F_2^2$. The bijection on vertices is
\[
          \phi:(z^\epsilon t(u),e)\longmapsto(u,\epsilon,e).
\]
Let $S=S^{-1}$ be a connection set and put $T=\phi(S)$.
For a noncentral quaternion $q$, we have $q^{-1}=zq$; hence
$(q,e)\in S$ if and only if $(zq,e)\in S$.
If two vertices have different $u$ coordinates, the quaternion part
of their difference has quotient coordinate $u+v$, while its sign
may be modified by the multiplication cocycle. Membership in $S$
ignores that sign. If they have the same $u$ coordinate, the
quaternion difference is exactly $z^{\epsilon+\eta}$.
Thus two vertices are adjacent precisely when the sum of their
images under $\phi$ lies in $T$, proving the claim.
\end{proof}

Payan's theorem says that a finite cubelike graph cannot have
chromatic number three \cite{Payan}; see also
\cite[Theorem~1.1]{KrebsSankar}. For arbitrary $E$, suppose a
graph in Lemma~\ref{lem:cubelike} is three-chromatic. It has a
proper three-coloring and an odd cycle. The finitely many $E$
coordinates on the cycle span a finite subgroup $E_0\leq E$.
Restricting both the coloring and the connection set to
$Q_8\times E_0$ produces a finite, nonbipartite, three-colorable
cubelike graph, contradicting Payan. Combining this with
Lemmas~\ref{lem:lift} and \ref{lem:structure} proves
Theorem~\ref{thm:main}.

\section{Computational verification and supporting files}

The accompanying packet includes the Python
script \texttt{verify\_small.py}. It checks the eight-vertex dihedral witness
and the quaternion-to-cubelike coordinate correspondence on inverse orbits
for $Q_8\times C_2^r$, $0\leq r\leq3$. It also enumerates all inverse-closed
connection sets for $r=0,1$. Of the $16$ sets for $Q_8$, $8$ give bipartite
graphs and $8$ give graphs that are not three-colorable. Of the $512$ sets for
$Q_8\times C_2$, the corresponding counts are $113$ and $399$. These finite
checks support the constructions; the proof of Theorem~\ref{thm:main} is the
argument above.

The paper, source files, and finite Lean witness are deposited at
\url{https://github.com/alexgocardinal/krebs-sankar-q53}
(initial source commit \texttt{ce84e11}). The repository is currently
private and will become readable to the public when released. The
\texttt{KrebsSankarQ53/Witness.lean} module compiles under Lean 4.29.0;
\texttt{Audit.lean} checks its concrete three-chromatic Cayley witness.
The packet does not contain a complete Lean formalization of
Theorem~\ref{thm:main}, and we make no claim of full machine verification.
The separate development
\cite{CayleyExpanded} proves a conditional section criterion. Its successful
compilation would not, by itself, certify every assertion of the present
manuscript.

\medskip\noindent
\textbf{Prior work.}
Goldstone and Weld characterized the nonabelian graphically abelian
groups as $Q_8\times E$ \cite{GoldstoneWeld}. A public Lean development
\cite{CayleyExpanded} states a related $D_8$-section criterion conditional
on three mathematical inputs, including a finite no-$D_8$ structural
classification, and contains the same concrete dihedral witness.
The present proof gives the structural argument directly for arbitrary
groups, without the finite classification or a local-finiteness theorem.
The connection sets in Theorem~\ref{thm:main} need not generate the
ambient group; the connected version is a separate question.

\clearpage
\begin{thebibliography}{9}
\bibitem{KrebsSankar}
M. Krebs and M. Sankar,
\emph{Abelian groups without $3$-chromatic Cayley graphs},
J. Combin. Theory Ser. B \textbf{179} (2026), 205--218.
\href{https://doi.org/10.1016/j.jctb.2026.03.004}{doi:10.1016/j.jctb.2026.03.004};
\href{https://arxiv.org/abs/2410.11028}{arXiv:2410.11028v2}.

\bibitem{Payan}
C. Payan, \emph{On the chromatic number of cube-like graphs},
Discrete Math. \textbf{103} (1992), 271--277.
\href{https://doi.org/10.1016/0012-365X(92)90319-B}{doi:10.1016/0012-365X(92)90319-B}.

\bibitem{GoldstoneWeld}
R. Goldstone and K. Weld, \emph{Graphically abelian groups},
Discrete Math. \textbf{310} (2010), 2806--2810.

\bibitem{CayleyExpanded}
\emph{CayleyExpanded: Lean 4 formalization of a section criterion
for three-chromatic Cayley graphs of exponent-four groups},
public repository, \url{https://github.com/tysunseven/CayleyExpanded}
(accessed October 2026).
\end{thebibliography}

\end{document}
